\appendix

\section{Reproduction}
\label{app:repro}

Each experiment is one script that writes one JSON file; the text quotes the JSON.
All scripts are pure \texttt{numpy}/\texttt{scipy}, run on a CPU, and the whole
sequence takes about 70 minutes.  The scripts are named after the order in which
they were run, not after the paper's sections.

\begin{table}[h]
\centering
\caption{Measurement scripts.  ``cell'' counts are the number of distinct
(matrix, precision, step count) measurements the script performs, which is the
number the text quotes.}
\label{tab:scripts}
\footnotesize
\setlength{\tabcolsep}{4pt}
\begin{tabular}{>{\raggedright\arraybackslash}p{5.6cm}
                >{\raggedright\arraybackslash}p{6.3cm}
                >{\raggedright\arraybackslash}p{3.4cm}}
\toprule
script & what it measures & used in \\
\midrule
\texttt{pilot1\_polar\_precision} & exact vs p-bit error of NS5/PE/LSQ5 designs & \S\ref{sec:exp} \\
\texttt{pilot1c\_precision\_scaling} & error against $K$ at fixed $p$ (flatness) & \S\ref{sec:exp} \\
\texttt{pilot4\_precision\_placement} & two-operand table, equal-cost allocation & \S\ref{sec:alloc} \\
\texttt{pilot4b\_k1\_separation} & the same table for $K=1,2,3,6$ (quintic) & \S\ref{sec:alloc} \\
\texttt{pilot4c\_k1\_highdeg} & $K=1$ with a degree-17 step ($E_1$ small) & \S\ref{sec:firstorder}, \S\ref{sec:alloc} \\
\texttt{pilot5\_lower\_bound} & single-step LP frontier, degree scaling, $\alpha$ & Appendix~\ref{app:refute} \\
\texttt{pilot5b\_scaling} & $S_{\mathrm{peak}}(\ell)$, $K^\star(\ell)$, plateau constant & \S\ref{sec:exp} \\
\texttt{pilot5c\_composition\_design} & non-monotone design path, $S_{\min}(\varepsilon,d)$ & Appendix~\ref{app:refute} \\
\texttt{pilot6\_s\_pert} & 192-cell competition of certificate definitions & \S\ref{sec:prelim}, \S\ref{sec:cert} \\
\texttt{pilot6b\_floor\_constant} & injected rounding size, $d$-scaling, $u$-linearity & \S\ref{sec:cert} \\
\texttt{pilot7\_pert\_aware\_design} & 37 designs measured, objective argmin, divergence census & \S\ref{sec:design} \\
\texttt{pilot7c\_table1} & certificate table over $p=4..24$, 90-cell stopping & \S\ref{sec:cert} \\
\texttt{pilot8\_muon\_mlp} & coupled Muon runs, frozen replay, rules, dose-response & \S\ref{sec:loop} \\
\texttt{pilot8b\_equilibrium} & subspace/value split, tail extension & \S\ref{sec:invariance} \\
\texttt{pilot10\_subspace\_theory} & invariance, first-order formula, $C(\ell)$, SVD-free bound & \S\ref{sec:cert} \\
\texttt{pilot10b\_first\_step\_schedule} & slot profile, 62-cell stopping, 48-cell schedule & \S\ref{sec:stopping}, \S\ref{sec:alloc} \\
\texttt{pilot11\_momentum\_subspace} & 500 real momentum matrices: two-term formula, $C\ell$ & \S\ref{sec:realmomentum} \\
\texttt{pilot11b\_application\_order} & fixed design and matrix, $K=1,2,3,4$ & \S\ref{sec:alloc} \\
\texttt{pilot12\_arithmetic\_variants} & three other kernels, self-calibration of $c$, plan cost & \S\ref{sec:variants} \\
\texttt{pilot13\_ell\_sensitivity} & design and stopping rule under $\hat\ell=\alpha\ell$, 252 decisions & \S\ref{sec:ellsens} \\
\texttt{pilot14\_real\_allocation} & graded schedule on 500 real momenta, 712 matrix-cells, $\ell$ sweep & \S\ref{sec:premium} \\
\texttt{pilot15\_fp8\_kernel} & E4M3/E5M2 with subnormals and saturation, 5 seeds, 4 dynamic ranges & \S\ref{sec:fp8} \\
\texttt{pilot16\_stats\_and\_lp} & bootstrap intervals, LP grid convergence, integer DP, held-out shapes & \S\ref{sec:loop}, \S\ref{sec:alloc} \\
\texttt{pilot17\_theory\_refinement} & family dependence of the constant, pair recursion, safety boundary & \S\ref{sec:limitations} \\
\texttt{pilot18\_c\_scaling\_and\_e5m2\_fix} & dimension normalisation of the floor constant, four-bit repair sweep & \S\ref{sec:loop}, \S\ref{sec:fp8} \\
\texttt{pilot18b\_bootstrap} & intervals and paired gains on the two above & \S\ref{sec:loop}, \S\ref{sec:fp8} \\
\texttt{pilot19\_large\_d\_and\_decomp} & dimension extrapolation to $d=512$; the constant's factors & \S\ref{sec:loop} \\
\texttt{pilot20\_analytic\_and\_d1024} & $2\sqrt2$ from the theorem-3 operator; $d=768,1024$; the practical rules & \S\ref{sec:loop}, \S\ref{sec:rules} \\
\texttt{pilot21\_loss\_phase\_and\_scheduler} & the kernel inside training (digits, diabetes, 11 seed-runs); step budget on 80 real momenta; online sensors & \S\ref{sec:realtrain}, \S\ref{sec:steps}, \S\ref{sec:online} \\
\texttt{pilot21b\_kmin} & per-matrix step counts on the same 80 momenta & \S\ref{sec:steps} \\
\texttt{pilot22\_participation\_factor} & the profile factor $g$ on 180 cells; the block-operator closure; per-injection coefficients & \S\ref{sec:profile} \\
\texttt{pilot23\_divergence\_mechanism} & the window and the escape; the four counterfactuals; the loop gain & \S\ref{sec:escape} \\
\texttt{pilot24\_large\_d\_multiplicity} & the multiplicity to $d=512$; per-injection quartiles; the ratchet's bound & \S\ref{sec:profile}, Appendix~\ref{app:bound} \\
\texttt{pilot25\_long\_census} & 1,200 steps $\times$ 25 seeds: the hazard of a call, with intervals and tests & \S\ref{sec:realtrain}, \S\ref{sec:escape}, Table~\ref{tab:census} \\
\texttt{pilot25b\_seed\_level} & the same runs read per seed: run rates, seed bootstrap, seed-level tests & \S\ref{sec:escape}, Table~\ref{tab:seedlevel} \\
\texttt{make\_figure\_data} & every plotted coordinate, read from the JSON files above & Figs.~\ref{fig:ratchet}--\ref{fig:pareto} \\
\bottomrule
\end{tabular}
\end{table}

\paragraph{Determinism.}
The momentum sequence used in \S\ref{sec:realmomentum} is reproduced bit for bit
by rerunning the exact-arithmetic trajectory with the same seed; the replay was
checked against the sequence of the run that produced \S\ref{sec:loop} (median
dynamic range $0.1499$, best validation loss $0.6534$, accuracy $0.762$, all
identical).  The matrix of \S\ref{sec:alloc} is reproduced by seeding
\texttt{numpy} as in the single-step script; its exact error $E_1$ matches to
seven digits.

\paragraph{Design generator.}
The coefficient generator (a Newton iteration on the fifth-degree equioscillation point)
returns non-finite coefficients at five of three hundred grid values of $\ell$, from
$3.2\times10^{-2}$ to $0.46$.  The held-out grid that would have used $3\times10^{-2}$
tests the design and raises instead of running the kernel with a non-finite one, and uses
$5\times10^{-2}$ there.

\paragraph{Archive.}
The scripts, their JSON results and logs, and the validation script are distributed as
\path{iclr-2028-artifact.tar.gz}; \texttt{MANIFEST.sha256} pins each experiment file
and hashes to \texttt{05ab9d2132ab4f46}\allowbreak\texttt{39dfb44366811d45}\allowbreak\texttt{5f228d2b2a4e38d3}\allowbreak\texttt{fd48a389544bd567}, so the numbers in the tables can be checked
against the code that produced them without re-running it.

\section{The Negative Results in Full}
\label{app:refute}

\paragraph{No lower bound of the form $S\ge c\,\varepsilon^{-\alpha}$.}
Three measurements, all from linear programs over odd polynomials of degree
$d\in\{3,5,7,9,13,17\}$ with the constraint $\max_{[\ell,1]}|P-1|=\varepsilon$ and
the objective the coefficient $\ell_1$ norm.

(i) At fixed error the minimal sensitivity decreases monotonically in the degree:
\[ S_{\min} = 1.178,\ 0.937,\ 0.648,\ 0.371,\ 0.114,\ 0.0363
   \quad \text{for } d = 3,5,7,9,13,17, \]
at $\ell=0.3$, $\varepsilon=0.6$ --- a factor $32$ across the
family.  So small $\varepsilon$ and small $S$ are simultaneously achievable.

(ii) The local exponent of $S_{\min}(\varepsilon)$ is not a constant, and it depends
on $\ell$ by two orders of magnitude:

\begin{table}[h]
\centering
\caption{Local exponent $\alpha$ of $S_{\min}\propto\varepsilon^{-\alpha}$ at three
dynamic ranges, from the linear program (log--log fit over the admissible range of
$\varepsilon$; fit residual $0.19$--$0.26$ in $\log$).}
\label{tab:alpha}
\begin{tabular}{rcccccc}
\toprule
$\ell$ & $d=3$ & $d=5$ & $d=7$ & $d=9$ & $d=13$ & $d=17$ \\
\midrule
$0.3$   & 1.43 & 0.81 & 0.56 & 0.45 & 0.32 & 0.22 \\
$0.1$   & 7.50 & 3.55 & 2.39 & 1.78 & 1.13 & 0.77 \\
$0.03$  & 21.4 & 13.1 & 7.11 & 5.02 & 4.60 & 3.13 \\
\bottomrule
\end{tabular}
\end{table}

(iii) Along a design path at fixed $(\ell,K,d)$ the error falls by $500\times$ while
the sensitivity is non-monotone:
\[ \varepsilon = 0.843,\ 0.782,\ 0.705,\ 0.612,\ 0.505,\ 0.386,\ 0.263,\ 0.150,\ 0.061,\ 0.021,\ 0.0016, \]
\[ S = 155.6,\ 215.0,\ 286.4,\ 366.7,\ 448.5,\ 516.8,\ 548.8,\ 517.1,\ 404.3,\ 228.0,\ 60.4. \]
The local slope
$\mathrm{d}\log S/\mathrm{d}\log\varepsilon$ runs from $-4.26$ to $+0.53$.

\paragraph{Why.}
The work of the lower bound would have to happen in the boundary layer
$[0,\ell]$, where the polynomial has to leave its value at the origin; the
sensitivity is measured on $[\ell,1]$ only, and raising the degree thins the
boundary layer without paying on the interval.  At $\ell=10^{-3}$ a single step
cannot achieve an error below $0.873$ even at degree $65$, so the regime the paper
works in is one where iteration is mandatory and no single-step bound can be tight.

\paragraph{Two operands.}
Table~\ref{tab:k1} is the two-operand question across the state and the intermediate.
The whole column of $B$ from 6 to 24 bits is worth $1.12\times$ at $p_X=8$ and
$71.7\times$ at $p_X=16$: the bottleneck is whichever operand is coarser, and at equal
total bits $(16,16)$ gives $\e{2.14}{-3}$ against $\e{2.90}{-2}$ for $(24,8)$.  The
operands are interchangeable to $1\%$ at $K=1$ and separate by $2.5\times$ at $K=2$, and
the quintic keeps that separation at $K=6$.  On the real spectra it is larger: $(24,8)$
gives $\e{5.74}{-2}$ against $\e{2.44}{-4}$ for $(16,16)$, a factor $236$.

\paragraph{The design axis is empty.}
\label{sec:design}
The law invites one more use: minimize $J(K)$ over the coefficients.  Three measurements
of a $37$-design family (cushion, safety, first-step damping $\tau$; $d=64$) say no.
(i) \emph{No interior optimum}: the argmin over
$(\text{cushion},\text{safety},\tau,K)$ sits on the scan boundary ($K=12$, $\tau=2.2$)
in every cell, because a step costs arithmetic and $J$ does not count it.  (ii) \emph{The
last step's derivative is not a handle}: with $p'(1)=\lambda$ fixed, $J$ is minimized at
$\lambda^\star=-0.0353$, the shipped step's own value, while $S_{\mathrm{pert}}$ is
identical ($609.6$) for all $|\lambda|\le0.5$.  (iii) \emph{Measured, the shipped design
is already there}: at $\ell=10^{-3}$, $K=8$, bf16 the $35$ non-divergent designs span
$0.232$--$0.768$, the shipped point is $2.578\times10^{-1}$ ($1.112\times$ from the best),
and the objective's argmin is $1.58\times$ \emph{worse}.  The reason is in
\eqref{eq:firstorder}: the subspace gain is set by the whole spectrum while
$S_{\mathrm{pert}}$ is a scalar of the one-dimensional recurrence, so $c$ is not
design-independent --- over the 37 designs it ranges over $-0.08$--$0.25$ with
$\mathrm{Spearman}(S_{\mathrm{pert}},\mathrm{err})$ between $-0.84$ and $-0.91$.  Designs
with safety $1.0$ diverge at bf16 in six of the 37 cells, so the shipped $1.01$ is a
divergence margin (\S\ref{sec:limitations}).

\section{A Matrix-Free Bound, and Why No Spectral Bound Exists}
\label{app:bound}

\begin{proposition}[Computable upper bound]
\label{prop:bound}
For full-rank $A$, $B$,
$\norm{\polar(A)-\polar(B)}_2 \le 2\norm{A-B}_2/(\sigma_{\min}(A)+\sigma_{\min}(B))$.
With $A = q(\hat M,p)$, $B = \hat M$ and $\sigma_{\min}(A) \ge \ell - \norm{E}_2$,
\begin{equation}
\mathrm{sub}_{\mathrm{err}} \;\le\; \frac{2\norm{E}_2}{2\ell - \norm{E}_2},
\qquad E = q(\hat M,p) - \hat M ,
\label{eq:bound}
\end{equation}
where $\ell$ is the design's assumed dynamic range and $\norm{E}_2$ is obtained
with matrix--vector products only: power iteration on $E^\top E$ reproduces the
SVD value to a relative $1.2\times10^{-15}$ in all 24 bound cells.
\end{proposition}

Equation~\eqref{eq:bound} is valid in all 24 cells tested and loose by
$2.1$--$14.0\times$ (median $5.7\times$); the slack grows with the dynamic range
because $1/(2\ell)$ assumes the worst orientation while rounding has a fixed
participation.  Replacing $\norm{E}_2$ by the exactly computable $\norm{E}_F$ costs
another factor: the Frobenius form is valid but loose by $6.3$--$45\times$ (median
$24\times$).  A matrix-free guarantee therefore exists, at about one order of
magnitude in accuracy; the tight prediction of \eqref{eq:firstorder} is not
matrix-free.

\begin{proposition}[Spectral information is not enough]
\label{prop:nospectral}
No function of the singular values of $X$ upper-bounds
$\mathrm{sub}_{\mathrm{err}}(X) = \norm{\polar(X)-\polar(M)}_2$.
\end{proposition}

Take a random orthogonal $Q$: its singular values are all $1$, so any spectral
functional takes the value it takes at $\polar(M)$, while
$\norm{\polar(Q)-\polar(M)}_2 = 2.000$ (measured: $\norm{QQ^\top-I}_2 =
1.6\times10^{-15}$).  The subspace is free once the spectrum is fixed.  On our grid
the two quantities co-vary ($\mathrm{Spearman}(\norm{XX^\top-I},
\mathrm{sub}_{\mathrm{err}}) = +0.83$, 32 cells), but only because both grow with
the dynamic range: a correlation is not a certificate.  This is why
\eqref{eq:cert} contains the orbit sensitivity of the \emph{design}, and why the
floor analysis contains $\norm{E}$, the size of the start error, rather than
anything measured from $X$ alone.

\paragraph{The ratchet, and the channel it certifies.}
Definition~\ref{def:ratchet} is a maximum, and it is worth stating exactly what it
bounds --- and what it does not.  Write the $p$-bit scalar orbit as
$\hat x_t = p_t(\hat x_{t-1}) + \delta_t$ with
$|\delta_t|\le\rho u$ --- one cast and, per step, the rounding of $B$ and of the state,
$2K+1$ injections in all --- and let $x_t$ be the exact orbit it shadows.

\begin{proposition}[What the ratchet bounds, and what it is for]
\label{prop:ratchet}
Let $p_l\in C^1(W)$ on a window $W$ containing the orbit of every $\sigma\in[\ell,1]$ and
its perturbations, $\rho=\max_W|x|$, and let $S(K)=\max_\sigma S_K(\sigma)$ with the empty
product --- the rounding of the output itself --- taken as $1$.  Then, for $u\to0$,
\begin{equation}
|x_K-\hat x_K| \;\le\; (2K+1)\,\rho\,u\,\max\{1,S(K)\} \;+\; O(u^2),
\label{eq:ratchet}
\end{equation}
and with the off-diagonal gain $\gamma_l(x,y)=a_l+b_l(x^2+y^2)+c_l(x^4+x^2y^2+y^4)$ in
place of $|p_l'|$ in the same maximum, call it $G(K)$, the subspace error obeys the same
inequality.  Neither inequality needs the ratchet: $S_{\mathrm{pert}}(K)\ge\max\{1,S(K)\}$,
so substituting it can only loosen \eqref{eq:ratchet}.  What the ratchet does is index the
law across $K$: $\max\{1,S(K)\}$ is not monotone in $K$, and the running maximum over
prefixes (Definition~\ref{def:ratchet}) is.
\end{proposition}

\begin{proof}
$\hat x_t = \phi_t(\hat x_{t-1})+\delta_t$ with $\phi_t$ differentiable on $W$, so the
first-order response to the injection at step $j$ is
$\delta_j\prod_{l>j}p_l'(x_{l-1})$.  For $j<K$ the product is
$\prod_{l=j+1}^{K}|p_l'|$, which is one of the products $S_K(\sigma)$ maximises over its
injection point --- $j'=j+1$ is among the points it maximises over --- and hence at most
$S(K)$; for the last injection, $j=K$, the product is empty and equals $1$.  Summing the
moduli over the $2K+1$ injections gives \eqref{eq:ratchet}.  The $O(u^2)$ term is the
second-order remainder of the composition, once per injection, and it can be given its own
coefficient: with $M=\max_W|p_l''|$ and $B=\max_W|x|$, the recursion
$e_t=p_t'(x_{t-1})e_{t-1}+\delta_t+\tfrac12p_t''(\xi_t)e_{t-1}^2$ with
$|e_{t-1}|\le(2t-1)BuS_{\mathrm{pert}}(t-1)$ propagates the quadratic parts through the
same prefix products, so
\begin{equation}
|x_K-\hat x_K| \le (2K{+}1)BuS_{\mathrm{pert}}(K)
+ \tfrac12 M B^2 u^2 S_{\mathrm{pert}}(K)\sum_{t=1}^{K}(2t-1)^2 S_{\mathrm{pert}}(t-1)^2
+ o(u^2).
\label{eq:second}
\end{equation}
At the shipped point ($\ell=10^{-3}$, $K=7$, $u=2^{-8}$, $B=1$, $M=2.0\cdot10^3$ over the
visited window) the quadratic coefficient exceeds the first-order term by orders of
magnitude: \eqref{eq:second} is a $u\to0$ statement, and at four bits what certifies
\eqref{eq:ratchet} is the measurement (24 of 24 cells, $16$--$1048\times$ slack), not the
remainder bound.  We claim it as a first-order envelope at four bits.  The pair channel is
the same argument with $\gamma_l$ in place of $|p_l'|$: an injection's off-diagonal
component is multiplied, step by step, by $\gamma_l$ at the pair's own orbit, so the
triangle inequality and the definition of $G(K)$ give it.
\end{proof}

\paragraph{Why the scalar channel collapses and the subspace one does not.}  The derivation
above uses $S(K)$ and an earlier form of this proposition concluded that
$S_{\mathrm{pert}}(K)$ could not be replaced by it, which was a contradiction: the two
claims answer different questions, and they are now separated.  The terminal value is the
valid bound; the ratchet is the predictor; and why they differ so far is measured rather
than argued.  At $\ell=10^{-3}$ the shipped design's prefixes give $S(t)=609$, $60$,
$2.5\cdot10^{-3}$ for $t=5,6,7$ --- the Newton-type last step has $p'(1)=0$ and multiplies
every suffix product --- so $\max\{1,S(K)\}$ is $1$ at $K=7$ and puts the two-term law at
$5.5\cdot10^{-4}$ where $2.58\cdot10^{-1}$ is measured, a factor $470$, while the ratchet
stays at $610$ and puts it at $0.34$.  The pair envelope over the same three counts is
$737$, $735$, $735$: the gain of the last step at the top of the window is
$\gamma(1,1)=1/2$ where $p'(1)=0$, and the subspace channel therefore does not inherit the
collapse of the scalar one.  That is also why the scalar $S$ enters the law as a calibrated
surrogate rather than as the certificate, and why the certificate of the floor is the pair
envelope, whose constant has to be measured ($0.011$--$0.912$) instead of read off.

Two measurements test it (\texttt{pilot24}).  The certified bound holds in all 24 cells
tested ($d=64$, $p=8$, $K=6,7$, four families, three $\ell$) with a slack of
$16$--$1048\times$ (median $41\times$), computed with $\rho=1$: the inequality is true and
loose by about the ratio of $2K+1$ to the calibrated $c$, which is what sign cancellation
does (\S\ref{sec:profile}).  The pair envelope is not a small correction to the scalar one:
$G(K)/S_{\mathrm{pert}}(K)$ is $0.62$--$1.21$ (below one in 16 of 24 cells), and
the constant it leaves, $\mathrm{err}/(u\,G(K))$, is $0.011$--$0.912$
(median $0.44$) --- below one in every cell.  So the subspace channel needs no constant
above one, and the $0.142$ the text reports is the scalar surrogate of that envelope.  The
two constants are unstable and tight for the same reason.  The scalar inequality replaces
each signed response by its modulus, so the $2K+1$ injections add at full size and the
measured error sits $41\times$ (median) below the bound only through a cancellation that
belongs to the data, not to the design --- which is why what the modulus leaves over moves
over two orders of magnitude.  The pair channel keeps each pair's off-diagonal gain and,
unlike the modulus, does not let the diagonal and off-diagonal parts of an injection be
traded against each other from step to step; that is what makes its envelope hold at one in
every cell (\S\ref{sec:profile}), and its price is a constant that has to be measured
rather than read off $S_{\mathrm{pert}}$.

What no condition on the spectrum or on the coefficients can do is make
$S_{\mathrm{pert}}$ a bound for the \emph{measured} floor.  The condition that would is
$\gamma_l(x,y)\le\max_{u\in[\min(x,y),\max(x,y)]}|p_l'(u)|$ on every reachable pair, and it
fails wherever a step has an interior zero of $p'$: a Newton step has $p'(1)=0$ and
$\gamma(1,1)=1/2$, and the ratio is unbounded as the pair approaches that point.  Over the
reachable pairs and steps the worst ratio is $1.8\times10^3$ at $\ell=10^{-3}$ and
$9\times10^9$ at $\ell=10^{-1}$, so the domination holds in none of the nine designs
tested, and the shipped sequence's own last step at $\ell=10^{-3}$, $K=7$ is
$(1.875,-1.250,0.375)$: $p'(1)=0$, $\gamma(1,1)=0.500$.  That is why the text reports
$S_{\mathrm{pert}}$ with a calibrated constant and names the pair channel as what the
multiplicity counts, instead of claiming a theorem for the floor it measures.

\section{Is the Kernel Model the Arithmetic?}
\label{sec:variants}

A reviewer is entitled to ask whether $c=0.142$ is a constant of our kernel model rather
than of the arithmetic.  We re-ran the shipped design under three other kernels on 12
cells each ($\ell\in\{10^{-2},10^{-3},10^{-4}\}$, $d\in\{64,128\}$, $p\in\{8,11,16\}$)
and report predicted/measured in Table~\ref{tab:variants}.  \emph{float32 accumulation}
gives $0.84$--$1.70$ (median $1.13$) and the same $2.629\times10^{-1}$ at the reference
cell as our kernel, so the wide-accumulator assumption is not what the constant measures;
\emph{eager rounding} raises the floor by up to $2.2\times$ (ratio $0.45$--$1.68$) and
\emph{stochastic rounding} gives $0.56$--$1.70$.  The constant moves by about a factor $2$
across kernels while the $u$-scale and the $K$-dependence do not.  A caller who does not
trust $0.142$ can fit it: $\hat c=0.109$ from 8 cells predicts a held-out grid to
$0.58$--$1.19$ (median $0.85$, CI $[0.64,1.12]$), for the cost of 8 runs; the plan it
enables takes $0.47$\,ms against $2.3$\,ms for the iteration.

\section{The Floor Constant across Families and Dimensions}
\label{app:constant}

\paragraph{The cast factor.}  Measured against the assumed $\rho=u/\sqrt3$ of the input
cast, the kernel's own cast rms is $1.472$ times larger, and for a log-uniform mantissa
$\sqrt{\mathbb{E}[m^{-2}]}=\sqrt{3/(2\ln2)}=1.471$ --- derived, and reproducing the
measured cast error to $0.6\%$ at every $d$ tested.  What remains is the response of an
incoherent perturbation: on the sixteen cells checked in this way the channel factor is
$4.53$, carried $3.8$ by the input cast and $2.3$ by the in-loop roundings in quadrature
($4.43$, $2\%$ low), so the 240-cell constant is $6.07=1.471\times4.13$.  Gaussian noise
of the same rms reproduces the cast channel to $1.00\times$ and a rank-one perturbation of
the same norm is $4\times$ larger, so the incoherence is real and the channel factor is
what stays measured.

\paragraph{Extrapolation.}  Extending the grid to $d=384$ and $512$ (four times the top of
the fit, 96 cells) leaves the constant at $6.48$, $+6.6\%$ with $88\%$ of cells inside
$\pm40\%$, and the trend fitted on $d\le256$ extrapolates to within $11\%$ of them.  Two
more dimension steps move nothing: over 32 cells at $d=768$ and $1024$ (two flat-start
families, two $\ell$, two $p$, two seeds) the median is $6.99$, $1.01\times$ the pooled
median of the same two families at $d=512$, every cell is inside $\pm40\%$ of $6.07$, and
the trend fitted on $d\le512$ extrapolates to $0.93$--$1.00$ of the measurement on all
$16$ predictions.

\paragraph{The participation, and its closure.}  Of the response factor $4.13$ behind
$6.07$, $2.83$ is derived.  On matrix units the first term of \eqref{eq:firstorder} is the
fixed operator $T(e_{ij})=(\sigma_i+\sigma_j)^{-1}(e_{ij}-e_{ji})$, checked against finite
differences over all $992$ units of a $d=32$ cell to a residual of $10^{-6}$, so its
Hilbert--Schmidt norm is $2(\norm{W}_F^2-\sum_iW_{ii}^2)^{1/2}$ without a fitted constant;
an exchangeable perturbation is then answered by $2\sqrt2\,\rho\norm{W}_F/\sqrt d$ --- a
Wigner factor $2$ and a skew factor $\sqrt2$ --- and a flat spectrum realises it ($2.81$
against $2\sqrt2=2.83$ at $d=1024$).  A real spectrum breaks exchangeability, and the
factor that survives, $g$, is measured over 180 cells at $1.31$ (range $0.98$--$2.10$).
It closes as the edge of a coarse-grained variance-profile operator: with $V=2D^2$ and $K$
contiguous blocks of equal index count, $B_{kl}=\overline{V}_{kl}\sqrt{m_km_l}$ and the
prediction $2\sqrt{\lambda_{\max}(B)}$ are within $1.2\%$ of the measurement at the
median for $K=4$ and inside $\pm10\%$ in $97\%$ of 120 cells; the profile's
$\lambda_{\max}$ is $11$--$23\times$ too large and a row degree $76$--$1657\times$, both
of which are this operator with the normalisation misplaced.  The loop is not a second
factor: the one-step formula and the full $K$-step propagation of the same error agree to
$0.95$--$1.07$ (median $0.998$).

\paragraph{What the loop does add, counted.}  The kernel injects one cast and two errors
per step, and their first-order responses add in quadrature to $0.92$--$1.12$ of the
measured floor over 96 cells, with $84\%$ of the total on the input cast.  Giving every
injection the one-step law of the state it enters, however, over-predicts the floor by
$1.27$--$2.36\times$: the per-injection coefficient divided by $2\sqrt2\,g(s_t)$ spans
$0.003$--$2.6$ (median $0.57$), because a rounding injected into an intermediate is not a
rounding of the state.  So $4.13=2.828\times1.46$, with $1.46=g\,m$: $g=1.31$ on the 180 profile cells,
and $m=1.11$ the counted multiplicity, both computable and neither fitted.

\paragraph{The coefficient, in quartiles.}  A range of $0.003$--$2.6$ hides where the mass
is; Table~\ref{tab:w13} splits it.  The input cast obeys the one-step law almost exactly
(median $0.990$, IQR $[0.944,1.011]$ over $16$ injections), the in-loop state roundings are
within $11\%$ of it at the median, and the intermediate's own rounding is the one that
carries almost no first-order response (median $0.054$): the spread the text quotes is one
channel, not noise, which is exactly why the multiplicity is counted rather than modelled
by a per-state law.  The median is stable in $d$ and in $\ell$ (Table~\ref{tab:w13}), and
the coefficient itself does not drift with $d$ where $g$ is left out: over $1248$
injections of the 96 large-$d$ cells the raw $\kappa$ has median $2.11$, $2.21$, $2.25$,
$2.29$ at $d=64,128,256,512$ (a flat spectrum gives $2\sqrt2=2.83$).

\begin{table}[h]
\centering
\caption{The per-injection coefficient of \S\ref{sec:profile} as quartiles
(\texttt{pilot22}, $208$ injections, $p=8$, $K=6$, $d=64,96$): the response of one rounding
divided by $2\sqrt2\,g(s_t)\,\mathrm{rms}(E)\,\mathrm{scale}_t$, so that $1$ means the
one-step law of the state the rounding enters.  ``$X$'' is the rounding of the state,
``$B$'' the rounding of the intermediate; $\ell=10^{-1}$ and $10^{-2}$ are the two windows
of the grid.}
\label{tab:w13}
\begin{tabular}{lrrrrr}
\toprule
group & min & Q1 & median & Q3 & max \\
\midrule
all injections      & $0.003$ & $0.103$ & $0.569$ & $0.931$ & $2.64$ \\
input cast          & $0.878$ & $0.944$ & $0.990$ & $1.011$ & $1.10$ \\
in-loop, $X$        & $0.469$ & $0.715$ & $0.890$ & $0.959$ & $2.64$ \\
in-loop, $B$        & $0.003$ & $0.012$ & $0.054$ & $0.236$ & $2.21$ \\
$\ell=10^{-1}$      & $0.003$ & $0.038$ & $0.696$ & $0.902$ & $1.10$ \\
$\ell=10^{-2}$      & $0.010$ & $0.149$ & $0.539$ & $0.951$ & $2.64$ \\
logspace            & $0.003$ & $0.149$ & $0.536$ & $0.900$ & $1.05$ \\
two-plateau         & $0.003$ & $0.103$ & $0.720$ & $0.958$ & $2.64$ \\
gap                 & $0.003$ & $0.103$ & $0.541$ & $0.910$ & $1.04$ \\
power               & $0.003$ & $0.050$ & $0.533$ & $0.896$ & $1.10$ \\
\bottomrule
\end{tabular}
\end{table}

\paragraph{The multiplicity at large $d$.}  The quadrature of the 13 responses was
measured on 16 cells at $d=64,96$; does it survive the dimensions a real layer has?  It
does: over 96 cells (four families, two $\ell$, three seeds, $d=64,128,256,512$) the ratio
of the quadrature to the measured floor is $0.92$--$1.12$ (median $1.03$) with $97\%$ of
cells inside $\pm10\%$, and at $d=512$ alone $0.997$--$1.080$ (median $1.035$, 24 cells;
Table~\ref{tab:w13} and \texttt{pilot24}).  The responses themselves say why: the mean
$|\cos|$ between the 13 response matrices is $0.017$ at the median and falls with $d$
($0.023$, $0.015$, $0.012$, $0.011$ from $d=64$ to $512$), as an incoherent pair of
$d^2$-dimensional responses should, while the largest single cosine is $0.69$.  The cast carries
$0.85$--$0.86$ of the sum and $C_{\mathrm{total}}$ is flat in $d$ ($3.57$, $3.82$,
$3.67$, $3.80$ for the four dimensions).

\paragraph{The escape.}  The failing runs of \S\ref{sec:escape} repeat one trace.  At
diabetes seed 0, step 24, layer 0 the exact tops are
$0.70,\ 1.87,\ 1.65,\ 1.19,\ 1.01,\ 1.00$ and the E5M2 tops
$0.71,\ 1.93,\ 1.96,\ 3.17,\ 78,\ 1.0\cdot10^{9},\ 4.9\cdot10^{44}$ against certified
upper edges $1.00,\ 1.92,\ 1.71,\ 1.29,\ 1.02,\ 1.00$; the same shape appears at diabetes
seed 2 (step 4) and at digits seeds 0 and 1 (steps 327 and 381).  The step at which a run
escapes is a property of this implementation and this seed --- the same arithmetic escapes
at step 4 for another seed of the same task --- so what repeats is the trace, not its
index.  Beyond a window's edge
the quintic's gain is $2.2\times$ at $1.5$ times the edge and $9.3\times$ at twice, and
the largest start $1+\delta$ whose $p$-bit orbit over the six steps stays below $10$ is
$\delta=0.0075$ at $s=1.01$ for four bits, $0.048$ at $1.035$ and $0.063$ at $1.05$
($0.010$ and $0.035$ at five bits).  With the state rescaled into its window at every step
all 11 runs finish, and all 50 of the long runs below; with the learning rate cut by $10\%$
seven of 11 still fail.  The loop
gain $G=\lambda_{\max}(\mathrm{diag}(\mathrm{step})^{1/2}H\,\mathrm{diag}(\mathrm{step})^{1/2})$,
estimated by power iteration on Hessian-vector products (central differences of the
analytic gradient) every 25 steps, exceeds one on no trajectory (median $\le0.26$; maximum
$0.83$ in five of six; one $7.3$ spike in the sixth).

\section{The Census at 1,200 Steps}
\label{app:census}

\begin{table}[h]
\centering
\caption{The long census (\texttt{pilot25}): 1,200 training steps $\times$ 25 seeds per cell,
the hazard of one Muon call with a 95\% Wilson interval, and the calls at risk.  A call is
one application of the composite to one hidden weight matrix; a run that has not diverged
at 1,200 steps contributes all of its calls and no event (right censoring).  ``clamp'' is
E5M2 at $s=1.01$ with the state rescaled into its window whenever a step would leave it
(\S\ref{sec:escape}).}
\label{tab:census}
\begin{tabular}{llccccc}
\toprule
task & condition & div/runs & calls & hazard & 95\% CI & $1$ in \\
\midrule
\texttt{digits} & E5M2, $s=1.01$   & $25/25$ & $40\,002$  & $6.25\cdot10^{-4}$ & $[4.23,\,9.22]\cdot10^{-4}$ & $1\,600$ \\
        & E5M2, $s=1.035$ & $17/25$ & $99\,822$  & $1.70\cdot10^{-4}$ & $[1.06,\,2.73]\cdot10^{-4}$ & $5\,872$ \\
        & clamp           & $0/25$  & $120\,000$ & $0$ & $[0,\,3.20]\cdot10^{-5}$ & --- \\
\texttt{diabetes} & E5M2, $s=1.01$ & $25/25$ & $2\,730$ & $9.16\cdot10^{-3}$ & $[6.21,\,13.48]\cdot10^{-3}$ & $109$ \\
        & E5M2, $s=1.035$ & $0/25$  & $90\,000$  & $0$ & $[0,\,4.27]\cdot10^{-5}$ & --- \\
        & clamp           & $0/25$  & $90\,000$  & $0$ & $[0,\,4.27]\cdot10^{-5}$ & --- \\
\bottomrule
\end{tabular}
\end{table}

\begin{table}[h]
\centering
\caption{The same census at the level of the seed (\texttt{pilot25b}): 25 runs per cell, the
run rate with an exact Clopper--Pearson interval, the pooled call hazard with a bootstrap
that resamples \emph{seeds}, and the design effect --- that bootstrap's standard error
divided by the binomial standard error the Wilson interval uses.}
\label{tab:seedlevel}
\begin{tabular}{llccccc}
\toprule
task & condition & div/runs & run rate [95\%] & hazard & seed-boot 95\% & DEFF \\
\midrule
\texttt{digits} & $s=1.01$  & $25/25$ & $1.00\ [0.86,1.00]$ & $6.25\cdot10^{-4}$ & $[5.48,7.23]\cdot10^{-4}$ & $0.36$ \\
                & $s=1.035$ & $17/25$ & $0.68\ [0.46,0.85]$ & $1.70\cdot10^{-4}$ & $[1.14,2.26]\cdot10^{-4}$ & $0.69$ \\
                & clamp     & $0/25$  & $0\ [0,0.14]$      & $0$ & $0$ & --- \\
\texttt{diabetes} & $s=1.01$  & $25/25$ & $1.00\ [0.86,1.00]$ & $9.16\cdot10^{-3}$ & $[6.64,12.8]\cdot10^{-3}$ & $0.89$ \\
                & $s=1.035$ & $0/25$  & $0\ [0,0.14]$      & $0$ & $0$ & --- \\
                & clamp     & $0/25$  & $0\ [0,0.14]$      & $0$ & $0$ & --- \\
\bottomrule
\end{tabular}
\end{table}

\paragraph{The seed as the unit of inference.}  Calls within a run share weights and
momentum, so the $25$ seeds, not the calls, are the independent repetitions.  Resampling
seeds and pooling their calls gives the last two columns of Table~\ref{tab:seedlevel}: the
pooled hazard moves little ($6.25\cdot10^{-4}$ against the Wilson $[4.23,9.22]\cdot10^{-4}$
at \texttt{digits}, $1.01$) because exposure varies between runs more than the events do,
and the bootstrap's standard error is $0.36$--$0.89$ of the binomial one --- the effective
sample size is the run, so the call-level interval is not anti-conservative, it is the
wrong experiment.  On counts the conditions do separate at seed level: \texttt{digits}
$25/25$ against $17/25$ at $1.035$ (Fisher $p=4.0\cdot10^{-3}$, run-rate difference
$[0.16,0.52]$, hazard ratio $3.7$ with a seed bootstrap $[2.7,5.5]$) and \texttt{diabetes}
$25/25$ against $0/25$ ($p=1.6\cdot10^{-14}$); between the tasks the counts are tied
($25/25$ each, $p=1$) and what differs is \emph{when} a run fails, median $396$ steps
against $35$ (permutation $p=2.5\cdot10^{-5}$, Mann--Whitney $1.8\cdot10^{-9}$).  The
call-level exact $p$ below is therefore an ordering of hazards, and the seed-level tests are
the inference.

\paragraph{What the run count cannot say.}  At this horizon the two tasks are
indistinguishable as runs --- 25 of 25 fail in both at $s=1.01$, so Fisher's exact test is
undefined --- and they are $14.7\times$ apart per call: $42\,732$ calls carry 50 events, and
the exact conditional binomial test on that split gives $p=3.5\cdot10^{-17}$.  The task is
not a single scale factor either: at $1.035$ \texttt{digits} still fails at $1$ in
$5\,872$ calls (CI $1$ in $9\,400$ to $1$ in $3\,670$) while \texttt{diabetes} shows no
failure in $90\,000$ calls (upper bound $1$ in $23\,400$), so the ordering of the two tasks
is not preserved by the safety --- we report that as measured and offer no mechanism for
it.  What the safety and the clamp do is bounded the same way: $1.035$ lowers the
\texttt{digits} hazard $3.7\times$ and \texttt{diabetes} below $1$ in $23\,400$, the clamp
below $1$ in $31\,200$ on \texttt{digits} and $1$ in $23\,400$ on \texttt{diabetes} over
$210\,000$ calls in all.

\paragraph{The shape of the hazard in time.}  It is not stationary.  On \texttt{digits} at
$s=1.01$ the events per 200-step block are $3$, $10$, $11$, $0$, $1$ over $19\,273$,
$15\,740$, $4\,060$, $800$ and $129$ calls ($0.016\%$, $0.064\%$, $0.271\%$, $0.775\%$ in
the blocks that still have runs at risk), and at $1.035$ the first $400$ steps carry none
of the $17$ events.  On \texttt{diabetes} every event is in the first 200 steps (median
$35$ steps, mean $36.8$): that task fails early or not at all, which is why its $400$-step
census already caught it and why its $1.035$ row is empty.  A run that has survived the
first few hundred steps is therefore not safe, only slower --- the reading a per-run count
of $1,200$ steps cannot give and a per-call rate can.

\paragraph{The prefix check.}  Every 1,200-step run reproduces the $400$-step census of
\texttt{pilot23} on its own prefix: 24 of 24 comparable runs agree (five \texttt{digits}
seeds and three \texttt{diabetes} seeds, three conditions each), which is what makes the
two censuses one experiment and not two.  The script writes no wall-clock field, so the
JSON is a function of the seeds and reruns byte for byte.

\paragraph{Two caveats on the estimate.}  The Wilson interval treats each call as an
independent Bernoulli trial; within a run the calls are correlated through the weights, so
the interval describes the exposure rather than upper-bounding a dependent process, and
Table~\ref{tab:seedlevel} prices that gap with the seed bootstrap.  The
clamp is not free: the implementation costs a four-iteration power estimate per step, which
this paper does not price --- doing so at scale is the one measurement we would add to it.

\section{Additional Tables}
\label{app:tables}

Table~\ref{tab:firstorder} is the square-grid test of
Theorem~\ref{thm:firstorder}; Table~\ref{tab:twochannel} is the two-operand precision
sweep of \S\ref{sec:alloc}; Tables~\ref{tab:variants}--\ref{tab:ci} are the kernel
comparison of \S\ref{sec:variants}, the scale laws, the crossing rule, the two-operand
step-count comparison, the graded schedule against the premium, the fp8 comparison of
\S\ref{sec:fp8}, and the bootstrap intervals; Table~\ref{tab:dsinv} normalises the floor constant by
$\sqrt d$, and Table~\ref{tab:e5m2} is the four-bit repair sweep of \S\ref{sec:fp8}.
Table~\ref{tab:profile} is the per-slot gain profile behind \S\ref{sec:alloc},
Table~\ref{tab:ell} the wrong-dynamic-range decisions behind \S\ref{sec:ellsens},
Table~\ref{tab:realtrain} the seven training conditions of \S\ref{sec:realtrain},
Table~\ref{tab:steps} the step budget of \S\ref{sec:steps}, and
Table~\ref{tab:rules} collects the four rules of \S\ref{sec:rules}.

\paragraph{Scale laws.}  Table~\ref{tab:scaling} gives $S_{\mathrm{peak}}(\ell)$ and the
plateau constant; beside them, measured $\norm{\delta X}_2/u = 1.26$--$1.86$ and
$\norm{\delta B}_2/u = 0.96$--$5.5$ carry no $\sqrt d$ component, while $\mathrm{err}/u$
grows $1:1.07:1.16$ from $d=64$ to $128$ against $1:1.41:2.00$ for a bare $\sqrt d$ law
--- the $\sqrt d$ appears only inside the normalisation of the pair weight
(\S\ref{sec:loop}).

\begin{table}[h]
\centering
\caption{Per-slot gain profile $\beta_t$ of the subspace channel, $d=64$, $p=8$,
$K=6$.  The profile is independent of $K$: the input gain is $11.5$ and $59.3$
at $\ell=10^{-2},10^{-3}$ for every $K\in\{4,6,8,10\}$ tested, and the tail sits at
$1.1u$.}
\label{tab:profile}
\begin{tabular}{rccccccc}
\toprule
$\ell$ & input & step 0 & step 1 & step 2 & step 3 & step 4 & step 5 \\
\midrule
$10^{-2}$ & \textbf{11.5} & 6.2 & 3.1 & 1.6 & 1.1 & 1.1 & 1.1 \\
$10^{-3}$ & \textbf{59.3} & 29.0 & 9.8 & 6.4 & 2.1 & 1.3 & 1.1 \\
$10^{-4}$ & \textbf{512} & 222 & 55.5 & 18.0 & 8.7 & 4.1 & 1.9 \\
\bottomrule
\end{tabular}
\end{table}

\begin{table}[h]
\centering
\caption{The stopping rule with a wrong dynamic range, $\hat\ell=\alpha\ell$: 36 cells
and three targets each, the targets set relative to the true predicted floor.
``missed'' is a returned step count whose measured error exceeds $1.5\varepsilon$;
``refused'' counts cells where the rule refuses a target that some step count of the
true design meets within $1.5\varepsilon$ (each row also holds two cells where no step
count does).  ``understatement'' is $\max\mathrm{err}/J(K)$ on the believed window: how
far the caller's own prediction falls short of the measurement.}
\label{tab:ell}
\begin{tabular}{rcccccc}
\toprule
$\alpha$ & met & missed & refused & $\mathrm{err}/\varepsilon$ median & max & understatement \\
\midrule
$0.1$  & 18 & 0  & 16 & 0.30 & 0.99 & 1.0 \\
$0.25$ & 18 & 0  & 16 & 0.53 & 0.99 & 1.0 \\
$0.5$  & 26 & 0  & 8  & 0.32 & 0.98 & 1.0 \\
$1$    & 36 & 0  & 0  & 0.59 & 1.09 & 1.3 \\
$2$    & 20 & 16 & 0  & 1.25 & 59   & 101 \\
$4$    & 14 & 22 & 0  & 2.52 & 121  & 614 \\
$10$   & 12 & 24 & 0  & 3.11 & 244  & 1275 \\
\bottomrule
\end{tabular}
\end{table}

\begin{table}[h]
\centering
\caption{Theorem~\ref{thm:firstorder} against measurement at the input cast.
$C = \mathrm{sub}_{\mathrm{err}}/u$ measured; $C_{\mathrm{pred}}$ from
\eqref{eq:firstorder} with $E = q(\hat M,p)-\hat M$.  The cell marked $\dagger$
saturates at $\mathrm{sub}_{\mathrm{err}} = 2$.}
\label{tab:firstorder}
\begin{tabular}{rrrccc}
\toprule
$d$ & $\ell$ & $p$ & $C$ measured & $C_{\mathrm{pred}}$ & ratio \\
\midrule
32 & $10^{-1}$ & 8 & 2.24 & 2.24 & 1.00 \\
32 & $10^{-3}$ & 8 & 58.96 & 57.92 & 0.98 \\
32 & $3\cdot10^{-4}$ & 8 & $512.0^{\dagger}$ & 181.5 & $0.35^{\dagger}$ \\
64 & $10^{-1}$ & 8 & 2.14 & 2.14 & 1.00 \\
64 & $10^{-2}$ & 8 & 11.50 & 11.48 & 1.00 \\
64 & $10^{-3}$ & 8 & 59.30 & 58.04 & 0.98 \\
64 & $10^{-3}$ & 16 & 76.28 & 76.26 & 1.00 \\
64 & $3\cdot10^{-4}$ & 11 & 230.25 & 225.74 & 0.98 \\
128 & $10^{-3}$ & 8 & 77.51 & 77.05 & 0.99 \\
128 & $3\cdot10^{-4}$ & 8 & 207.79 & 238.66 & 1.15 \\
\bottomrule
\end{tabular}
\end{table}

\begin{table}[h]
\centering
\caption{Error $\|X_K - UV^\top\|_2$ against the storage precision of the state
$p_X$ (rows) and of the intermediate $p_B$ (columns); shipping design, $\ell=10^{-3}$,
$d=64$, $K=6$.  Each row and column is flat over a wide range and then falls
steeply --- the signature of a bottleneck that moves.}
\label{tab:twochannel}
\begin{tabular}{r|ccccccc}
\toprule
$p_X \backslash p_B$ & 6 & 8 & 10 & 11 & 13 & 16 & 24 \\
\midrule
6  & \e{1.52}{0} & \e{1.45}{0} & \e{1.39}{0} & \e{1.41}{0} & \e{1.40}{0} & \e{1.41}{0} & \e{1.40}{0} \\
8  & \e{3.53}{-1} & \e{3.38}{-1} & \e{3.07}{-1} & \e{3.00}{-1} & \e{3.19}{-1} & \e{3.17}{-1} & \e{3.15}{-1} \\
10 & \e{1.78}{-1} & \e{9.48}{-2} & \e{9.58}{-2} & \e{8.96}{-2} & \e{8.91}{-2} & \e{9.15}{-2} & \e{9.21}{-2} \\
11 & \e{1.55}{-1} & \e{4.26}{-2} & \e{3.48}{-2} & \e{3.73}{-2} & \e{3.43}{-2} & \e{3.49}{-2} & \e{3.51}{-2} \\
13 & \e{1.57}{-1} & \e{2.97}{-2} & \e{1.13}{-2} & \e{1.07}{-2} & \e{1.09}{-2} & \e{9.96}{-3} & \e{9.77}{-3} \\
16 & \e{1.52}{-1} & \e{2.87}{-2} & \e{8.33}{-3} & \e{4.83}{-3} & \e{2.37}{-3} & \e{2.14}{-3} & \e{2.12}{-3} \\
24 & \e{1.56}{-1} & \e{2.90}{-2} & \e{8.62}{-3} & \e{4.81}{-3} & \e{1.90}{-3} & \e{1.64}{-3} & \e{1.62}{-3} \\
\bottomrule
\end{tabular}
\end{table}

\begin{table}[h]
\centering
\caption{Predicted/measured error ratio of the shipped design under four kernels of
$p$-bit arithmetic (\S\ref{sec:variants}), at $K=6$, one matrix per
$(\ell,d)$.  ``eager'' rounds every matrix product, ``fp32 acc.'' accumulates in
float32, ``stochastic'' uses stochastic rounding (mean of three seeds).}
\label{tab:variants}
\begin{tabular}{rrrcccc}
\toprule
$\ell$ & $d$ & $p$ & kernel & eager & fp32 acc. & stochastic \\
\midrule
$10^{-2}$ & 64 & 8 & 0.83 & 0.45 & 0.84 & 0.62 \\
$10^{-2}$ & 64 & 11 & 0.94 & 0.47 & 0.92 & 0.61 \\
$10^{-2}$ & 64 & 16 & 0.82 & 0.50 & 0.86 & 0.56 \\
$10^{-3}$ & 64 & 8 & 1.29 & 0.65 & 1.29 & 0.75 \\
$10^{-3}$ & 64 & 11 & 1.29 & 0.76 & 1.29 & 0.82 \\
$10^{-3}$ & 64 & 16 & 1.28 & 1.04 & 1.24 & 1.21 \\
$10^{-3}$ & 128 & 8 & 0.98 & 0.58 & 0.98 & 0.65 \\
$10^{-3}$ & 128 & 11 & 0.94 & 0.64 & 0.93 & 0.72 \\
$10^{-3}$ & 128 & 16 & 1.33 & 0.93 & 1.32 & 1.07 \\
$10^{-4}$ & 64 & 8 & 1.70 & 1.68 & 1.70 & 1.70 \\
$10^{-4}$ & 64 & 11 & 1.45 & 1.14 & 1.43 & 1.12 \\
$10^{-4}$ & 64 & 16 & 1.03 & 1.02 & 1.03 & 1.02 \\
\bottomrule
\end{tabular}
\end{table}

\begin{table}[h]
\centering
\caption{Graded schedule against uniform precision at equal total bits, as a function
of the first-slot premium.  First row: the real momentum sample; rows 2--5: a
logspace spectrum of the same $128\times32$ shape, two seeds, $K\in\{6,8\}$.  At
$\ell=10^{-4}$ the design's own exact error $E_K$ dominates at the two smaller budgets,
so the gain collapses there and returns at the largest.}
\label{tab:premium}
\begin{tabular}{lcccc}
\toprule
spectrum & premium (bits) & $K{=}6$, 8 bits & $K{=}6$, 12 bits & $K{=}8$, 12 bits \\
\midrule
real $128\times32$ & 0.9 & 1.00 & 1.00 & 1.00 \\
logspace $10^{-1}$ & 1.6--1.7 & 1.09--1.19 & 1.07--1.16 & 1.05--1.13 \\
logspace $10^{-2}$ & 4.3--4.5 & 3.17--3.71 & 3.22--3.59 & 4.23--4.86 \\
logspace $10^{-3}$ & 7.3--7.5 & 6.38--8.10 & 9.13--9.79 & 10.2--15.4 \\
logspace $10^{-4}$ & 8.8--9.8 & 2.25--2.32 & 1.30--1.46 & 20.7--29.1 \\
\bottomrule
\end{tabular}
\end{table}

% scale laws table moved here from the body
\begin{table}[h]
\centering
\caption{Scale laws of the certificate and the floor (synthetic logspace spectra,
shipped design, $p=10$).  $S_{\mathrm{peak}}$ is the maximum of $S(K)$ over $K$,
with the local log--log slope against $\ell$; $C = \mathrm{err}/u$ is the plateau
constant, blank where the design has not converged by $K=8$.}
\label{tab:scaling}
\begin{tabular}{rcccccc}
\toprule
$\ell$ & $S_{\mathrm{peak}}$ & $K^\star$ & local slope & $C$ ($d=32$) & $C$ ($d=64$) & $C$ ($d=128$) \\
\midrule
$10^{-2}$ & 77.1 & 3 & $-1.67$ & 10.1 & 12.2 & 13.4 \\
$10^{-3}$ & 609 & 5 & $-1.88$ & 48.5 & 82.9 & 93.1 \\
$10^{-4}$ & $5.63\times10^{3}$ & 6 & $-1.81$ & -- & -- & -- \\
\bottomrule
\end{tabular}
\end{table}

\begin{table}[h]
\centering
\caption{The crossing rule \eqref{eq:cross} and the target rule \eqref{eq:target} on
frozen real momentum spectra (500 matrices, design window $\ell=10^{-3}$; $K$ swept to
10).  Costs are the measured error at the selected $K$ over the smallest measured error
in $K$.  At $p=24$ the crossing lands on the argmin while the target rule stops three
steps later, at the plateau.}
\label{tab:cross}
\begin{tabular}{rrc rrrc rr}
\toprule
$p$ & plateau (u) & $K_{\arg\min}$ & $K_{\mathrm{cross}}$ & cost & $K_{\mathrm{target}}$ & cost & $K_{\mathrm{target}}$, calibrated $c$ & cost \\
\midrule
6  & 15.8 & 6 & 2 & $4.01\times$ & 10 & $1.002\times$ & 6 & $1.000\times$ \\
8  & 16.0 & 6 & 4 & $9.11\times$ & 10 & $1.001\times$ & 6 & $1.000\times$ \\
11 & 16.0 & 7 & 5 & $16.7\times$ & 10 & $1.001\times$ & 6 & $1.006\times$ \\
16 & 16.0 & 7 & 6 & $6.68\times$ & 10 & $1.001\times$ & 7 & $1.000\times$ \\
24 & 16.1 & 7 & 7 & $1.00\times$ & 10 & $1.006\times$ & 7 & $1.000\times$ \\
\bottomrule
\end{tabular}
\end{table}

\begin{table}[h]
\centering
\caption{Equal-total-bit comparison of the two operands as a function of the step count.
Rows 1--2 hold the design and the matrix fixed and change only $K$; rows 3--4 are the
shipped quintic at $\ell=0.3$, $d=64$; the last row is the frozen real-spectrum replay at
$K=8$.  ``asymmetry'' is the ratio of the two off-diagonal entries.}
\label{tab:k1}
\footnotesize
\setlength{\tabcolsep}{3pt}
\begin{tabular}{>{\raggedright\arraybackslash}p{3.25cm}
                >{\raggedright\arraybackslash}p{2.0cm}
                ccccc}
\toprule
setting & $\ell$, $d$ & balanced & intermediate richer & state richer & ratios & asymmetry \\
\midrule
degree-17, $K=1$, 16 bits & 0.3, 64 & $(8,8)$: \e{1.13}{-2} & $(6,10)$: \e{3.29}{-2} & $(10,6)$: \e{3.33}{-2} & 2.90, 2.94 & $1.01\times$ \\
degree-17, $K=2$, 16 bits & 0.3, 64 & $(8,8)$: \e{1.03}{-2} & $(6,10)$: \e{3.38}{-2} & $(10,6)$: \e{8.32}{-2} & 3.29, 8.11 & $2.46\times$ \\
quintic, $K=2$, 16 bits & 0.3, 64 & $(8,8)$: \e{1.09}{-2} & $(6,10)$: \e{3.79}{-2} & $(10,6)$: \e{1.67}{-2} & 3.48, 1.53 & $2.27\times$ \\
quintic, $K=6$, 16 bits & 0.3, 64 & $(8,8)$: \e{1.21}{-2} & $(6,10)$: \e{4.04}{-2} & $(10,6)$: \e{1.78}{-2} & 3.33, 1.47 & $2.27\times$ \\
real spectra, $K=8$, 32 bits & 0.15, $128\times32$ & $(16,16)$: \e{2.44}{-4} & $(8,24)$: \e{2.82}{-2} & $(24,8)$: \e{5.74}{-2} & 116, 236 & $2.03\times$ \\
\bottomrule
\end{tabular}
\end{table}

\begin{table}[h]
\centering
\caption{Critical safety $s^\star$: the smallest damping at which no run of the
12-spectrum census diverges ($d\in\{48,64\}$, six seeds, $K=6$).  A dash marks a cell
that is already stable at $s=1.0$; the last column is the measured
$c=(s^\star-1)/2^{-p}$.}
\label{tab:safety}
\begin{tabular}{rccc c}
\toprule
$p$ & $\ell=10^{-1}$ & $\ell=10^{-2}$ & $\ell=10^{-3}$ & $c$ \\
\midrule
4  & 1.0028 & 1.0164 & 1.0313 & 0.04--0.50 \\
5  & -- & 1.0053 & 1.0060 & 0.17--0.19 \\
6  & -- & 1.0032 & 1.0047 & 0.20--0.30 \\
8  & -- & 1.0002 & 1.0005 & 0.05--0.12 \\
10 & -- & -- & 1.0001 & 0.08 \\
12 & -- & -- & 1.00001 & 0.04 \\
16 & -- & -- & -- & -- \\
\bottomrule
\end{tabular}
\end{table}

\begin{table}[h]
\centering
\caption{The floor constant with the pair operator normalised by $\sqrt d$
(\S\ref{sec:loop}): measured $(\mathrm{err}-E_K)/(u\,\hat v\norm{W}_F/\sqrt d)$ with
$W=[1/(\sigma_i+\sigma_j)]$ and $\hat v=\norm{\hat M}_F/(\sqrt3 d)$, median over two
$\ell$, two $p$ and three seeds per cell.  The last column is the held-out grid
($\ell\in\{5\cdot10^{-2},3\cdot10^{-3}\}$, $p\in\{6,16\}$, other seeds, $d\in\{48,128\}$).
The two columns after $d=256$ extend the same grid to four times its top, and at
$d=768$ and $1024$ the two flat-start families read $7.94$/$7.93$ and $6.11$/$6.17$.
Without the $\sqrt d$ these ratios fall with exponent $-0.42$ in $d$; weighting the
operator by $\sigma_i\sigma_j/(\sigma_i+\sigma_j)$ instead spreads $65\times$
($2.6\times$ worse out of sample) and normalising by its stable rank $136\times$
($2.9\times$).}
\label{tab:dsinv}
\begin{tabular}{lcccccccc}
\toprule
family & $d=32$ & $d=48$ & $d=64$ & $d=128$ & $d=256$ & $d=384$ & $d=512$ & held out \\
\midrule
logspace & 6.64 & 6.77 & 6.96 & 7.33 & 7.64 & 7.76 & 7.75 & 6.96 \\
two-plateau & 4.85 & 5.14 & 5.63 & 5.89 & 5.96 & 6.04 & 6.13 & 5.37 \\
gap & 5.44 & 5.56 & 5.85 & 6.11 & 6.27 & 6.40 & 6.40 & 5.96 \\
power & 5.30 & 6.00 & 6.50 & 7.35 & 7.25 & 7.49 & 7.74 & 6.26 \\
\bottomrule
\end{tabular}
\end{table}

\begin{table}[h]
\centering
\caption{Four-bit repairs at $K=6$ with the shipped design, 12-spectrum census
($d\in\{48,64\}$, six seeds).  ``blow-up'' counts runs whose state distance
$\norm{X_K-UV^\top}_2$ exceeds $5$; ``floor'' is the median subspace error of the census
(cap $2$); $s^\star$ is the smallest damping with no blow-up.  Each row names the
precision of the slots in order; ``cast'' is the input quantisation, E4M3 is $p=5$ and
bf16 $p=8$.}
\label{tab:e5m2}
\begin{tabular}{lcccccc}
\toprule
rule & blow-up & floor & floor & blow-up & floor & $s^\star$ \\
     & $\ell=10^{-2}$ & $10^{-2}$ & $10^{-1}$ & $\ell=10^{-3}$ & $10^{-3}$ & $10^{-3}$ \\
\midrule
$p=4$, $s=1.01$ & 4/12 & 0.877 & 0.217 & 12/12 & 2.00 & 1.0313 \\
$p=4$, $s=1.035$ & 0/12 & 0.841 & 0.221 & 0/12 & 2.00 & 1.0313 \\
cast bf16 & 5/12 & 0.419 & 0.165 & 10/12 & 2.00 & 1.0235 \\
cast, step 1 bf16 & 0/12 & 0.243 & 0.133 & 10/12 & 2.00 & 1.0221 \\
cast, steps 1--2 bf16 & 0/12 & 0.146 & 0.114 & 5/12 & 0.537 & 1.0238 \\
cast, step 1 E4M3 & 1/12 & 0.469 & 0.157 & 9/12 & 2.00 & 1.0219 \\
cast, steps 1--2 E4M3 & 0/12 & 0.416 & 0.140 & 9/12 & 2.00 & 1.0221 \\
E4M3 throughout & 0/12 & 0.411 & 0.107 & 0/12 & 1.92 & 1.0060 \\
bf16 throughout & 0/12 & 0.049 & 0.014 & 0/12 & 0.283 & -- \\
\bottomrule
\end{tabular}
\end{table}

\paragraph{The (format, step-count) grid.}  On \texttt{diabetes} (three seeds, 1200
steps) the training MSE at K=2 is $1.71\times10^{-3}$ for bf16 (24 bits, update error
$0.71$), $1.46\times10^{-3}$ for E4M3 (15 bits, $0.71$) and $2.14\times10^{-3}$ for E5M2
at safety $1.035$ (12 bits, $0.73$); at K=4 they are $1.58$, $1.60$ and
$1.62\times10^{-3}$, and no cell at $K\le4$ diverges.  Against the fp32 value $1.299\times10^{-3}$ the spread over the whole grid is
a factor $1.65$, while the update error spans two orders of magnitude.

\begin{table}[h]
\centering
\caption{The kernel inside training (\S\ref{sec:realtrain}).  Update error is the window
error at the last probe, median over seeds; ``div'' counts weight blow-up
($\max|W|>10^{3}$, reached before NaN).  Bits are the cast plus one stored width per
step.  All rows are $400$ steps; the $1{,}200$-step census is Table~\ref{tab:census}.
The E5M2 row at safety $1.01$ has no diabetes loss: all three runs stop.}
\label{tab:realtrain}
\begin{tabular}{lccccccc}
\toprule
condition & bits & upd.\ err & \texttt{digits} train & acc & div & \texttt{diabetes} train & div \\
\midrule
fp32              & 168 & $2.8\cdot10^{-14}$ & $5.84\cdot10^{-5}$ & $0.980$ & $0/5$ & $1.30\cdot10^{-3}$ & $0/3$ \\
bf16              & 56  & $1.6\cdot10^{-2}$  & $5.79\cdot10^{-5}$ & $0.977$ & $0/5$ & $1.38\cdot10^{-3}$ & $0/3$ \\
E4M3              & 35  & $1.2\cdot10^{-1}$  & $6.07\cdot10^{-5}$ & $0.980$ & $0/5$ & $1.46\cdot10^{-3}$ & $1/3$ \\
E5M2, $s=1.01$    & 28  & $2.4\cdot10^{-1}$  & $7.39\cdot10^{-5}$ & $0.977$ & $2/5$ & --- & $3/3$ \\
E5M2, $s=1.035$   & 28  & $2.5\cdot10^{-1}$  & $6.22\cdot10^{-5}$ & $0.980$ & $0/5$ & $1.57\cdot10^{-3}$ & $0/3$ \\
hybrid, E5M2 body & 40  & $1.4\cdot10^{-1}$  & $5.77\cdot10^{-5}$ & $0.980$ & $0/5$ & $1.44\cdot10^{-3}$ & $0/3$ \\
hybrid, E4M3 body & 44  & $7.3\cdot10^{-2}$  & $5.83\cdot10^{-5}$ & $0.977$ & $0/5$ & $1.42\cdot10^{-3}$ & $0/3$ \\
\bottomrule
\end{tabular}
\end{table}

\begin{table}[h]
\centering
\caption{Step count against precision on 80 real momentum matrices (\S\ref{sec:steps}),
design window $\ell_0=10^{-2}$; bits are the input cast plus one stored width per step,
and the error is the window error, median over the 80 matrices.  The last two rows are the
hybrid rule of \S\ref{sec:fp8} at $K=6$.}
\label{tab:steps}
\begin{tabular}{llccc}
\toprule
format & $K$ & bits & window error & \\
\midrule
fp32    & 6 & 168 & $2.8\cdot10^{-14}$ & \\
\midrule
E5M2    & 2 & 12 & $7.35\cdot10^{-1}$ & \\
        & 3 & 16 & $3.64\cdot10^{-1}$ & \\
        & 4 & 20 & $2.49\cdot10^{-1}$ & \\
        & 5 & 24 & $2.41\cdot10^{-1}$ & plateau \\
        & 6 & 28 & $2.38\cdot10^{-1}$ & \\
\midrule
E4M3    & 3 & 20 & $3.13\cdot10^{-1}$ & \\
        & 4 & 25 & $1.29\cdot10^{-1}$ & knee \\
        & 5 & 30 & $1.24\cdot10^{-1}$ & plateau \\
        & 6 & 35 & $1.25\cdot10^{-1}$ & \\
\midrule
bf16    & 2 & 24 & $7.17\cdot10^{-1}$ & \\
        & 3 & 32 & $2.96\cdot10^{-1}$ & \\
        & 4 & 40 & $2.43\cdot10^{-2}$ & knee \\
        & 5 & 48 & $1.60\cdot10^{-2}$ & plateau \\
        & 6 & 56 & $1.62\cdot10^{-2}$ & \\
\midrule
hybrid, E5M2 body & 6 & 40 & $1.45\cdot10^{-1}$ & \\
hybrid, E4M3 body & 6 & 44 & $7.36\cdot10^{-2}$ & \\
\bottomrule
\end{tabular}
\end{table}

\begin{table}[h]
\centering
\caption{The four rules of \S\ref{sec:rules}, the measurement behind each and its price.
``bits'' is the width of the input cast plus one stored width per step; a named dynamic
range means the logspace ladder at $d=64$, four seeds.}
\label{tab:rules}
\begin{tabular}{p{3.2cm}p{4.6cm}p{2.5cm}p{3.2cm}}
\toprule
decision & measurement & window error & cost \\
\midrule
stop at the plateau ($K=5$) & flat from $K=5$ in every format; sixth step $<1.5\%$; $78/80$ matrices prefer it & $1.60\cdot10^{-2}$ (bf16) & $+8$ bits, no gain \\
buy steps, not precision & $40$-bit hybrid $1.45\cdot10^{-1}$ against bf16 $K=4$ $2.43\cdot10^{-2}$; one downgraded slot $2.7$--$12.6\times$ & $6\times$ at equal bits & $4$ bits per slot saved \\
$p=4$ only with $s\ge1.035$ & blow-up $12/12\to0/12$ at $10^{-3}$; real data $6/11\to0/11$ at $400$ steps, $17/25$ at $1{,}200$ & cap $2.00$ in both & $+1.9\%$ at $10^{-1}$, $-4.1\%$ at $10^{-2}$ \\
$\hat\ell\le\ell$ (halve $\sigma_{\min}$) & error $\le1.07\times$, $\le1$ extra step & --- & one extra step \\
\bottomrule
\end{tabular}
\end{table}

\begin{table}[h]
\centering
\caption{fp8 against the same kernel with an unlimited exponent at the matched precision
(\S\ref{sec:fp8}; median of five seeds, $d=64$, $K=6$).  ``E4M3 scaled'' uses
per-tensor power-of-two scales; ``sub'' counts the stored singular directions of the
start that fall below the format's smallest normal.  $\dagger$ marks a cell whose error
sits at the cap $2$ that any pair of orthogonal factors obeys; at $\ell=1.3\times10^{-2}$
E5M2 saturates $8192$ stored values.}
\label{tab:fp8}
\begin{tabular}{rrccccc}
\toprule
$\ell$ & $p=5$ ideal & E4M3 & E4M3 scaled & $p=4$ ideal & E5M2 & sub \\
\midrule
$0.15$   & \e{9.86}{-2} & \e{9.89}{-2} & \e{1.03}{-1} & \e{1.85}{-1} & \e{1.84}{-1} & 0 \\
$0.1$    & \e{1.08}{-1} & \e{1.11}{-1} & \e{1.14}{-1} & \e{2.15}{-1} & \e{2.15}{-1} & 0 \\
$3.6\cdot10^{-2}$ & \e{1.69}{-1} & \e{1.71}{-1} & \e{1.71}{-1} & \e{3.48}{-1} & \e{3.58}{-1} & 0 \\
$1.3\cdot10^{-2}$ & \e{3.00}{-1} & \e{3.08}{-1} & \e{8.08}{-1} & $2.00^{\dagger}$ & $2.00^{\dagger}$ & 5 \\
$4.8\cdot10^{-3}$ & \e{7.21}{-1} & \e{7.67}{-1} & \e{1.40}{0} & $2.00^{\dagger}$ & $2.00^{\dagger}$ & 14 \\
\bottomrule
\end{tabular}
\end{table}

\begin{table}[h]
\centering
\caption{Bootstrap 95\% percentile intervals (10\,000 resamples, fixed seed) on every
ratio and constant the abstract and the introduction quote.  ``point'' is the statistic on the
observed sample.  The 500-cell intervals are narrow because the formula reproduces that
grid to three digits; the interval says the agreement is not carried by a few cells.}
\label{tab:ci}
\begin{tabular}{llcc}
\toprule
quantity & cells & point & interval \\
\midrule
certificate, err/pred median & 192 & 0.99 & $[0.98,1.00]$ \\
certificate, worst deviation & 192 & $1.60$ & $[1.59,1.60]$ \\
stopping rule, pred/meas median ($d=64$) & 90 & $1.11$ & $[1.07,1.14]$ \\
stopping rule, pred/meas median ($d=128$) & 90 & $1.10$ & $[1.08,1.13]$ \\
two-term formula, pred/meas median ($p=8$) & 500 & $1.0000$ & $[0.9999,1.0001]$ \\
self-calibrated prediction, held out & 9 & $0.85$ & $[0.64,1.12]$ \\
graded schedule, gain over uniform & 48 & $3.06$ & $[2.24,4.76]$ \\
certificate on held-out shapes and spectra & 48 & $0.42$ & $[0.37,0.48]$ \\
first-order formula on the same cells & 48 & $0.89$ & $[0.81,0.91]$ \\
dimension-normalised constant, fit & 240 & $6.07$ & $[5.98,6.14]$ \\
the same, held-out cells & 64 & $6.04$ & $[5.78,6.26]$ \\
\bottomrule
\end{tabular}
\end{table}

\paragraph{Discretisation checks.}
Two of the paper's quantities are produced by linear programs and one by integer
rounding, so each could be an artefact of its grid.  (i) The minimax single-step error at
$\ell=10^{-3}$ moves by $7.5\times10^{-4}$ (degree 5) and $2.2\times10^{-3}$ (degree 9)
between $n=50$ and $n=3200$ grid points, and by $2\times10^{-5}$ between $n=1600$ and
$n=3200$; degree 17 needs $n\ge400$ (a $5\%$ deviation at $n=50$, $1.9\times10^{-5}$ at
$n=3200$), and the values quoted in the tables use $n\ge400$.  (ii) The last-step family's coefficients at $p'(1)=\lambda$ move by
$5\times10^{-5}$ between $n=100$ and $n=1600$, and the argmin of $J$ over $\lambda$ is the
same grid point on $n\in\{100,400,1600\}$ with $J$ changing by $4\times10^{-6}$.
(iii) Against the exact integer optimum from a dynamic program over the $K+1$ slots, the
largest-remainder schedule of \eqref{eq:maxform} is exactly optimal in 44 of the 48
cells and within $1.47\times$ in the rest, so the quoted gain is not a rounding artefact.

\paragraph{A verification in fp8.}
Table~\ref{tab:fp8} also documents where the format leaves the law's domain: from
$\ell=1.3\times10^{-2}$ downward the smallest stored directions are subnormal (5 of 64
there, 14 at $4.8\times10^{-3}$), and below $\ell\approx6\times10^{-4}$ directions are
flushed and the error sits at the cap in every format, including the unlimited-exponent
one, because four significant bits are not enough for a dynamic range of $10^{3}$.
